\documentclass[12pt]{article}
\usepackage{amsmath}
\usepackage{amssymb}
\usepackage{geometry}
\usepackage{listings}
\usepackage{xcolor}
\usepackage{hyperref}
\usepackage{enumitem}
\usepackage{booktabs}

\geometry{margin=1in}

\title{MLOps for Federated Learning: Adapting Operations to Distributed Training}
\subtitle{How MLOps Principles Apply to Your DDoS Detection FL System}
\author{}
\date{}

\begin{document}

\maketitle

\section{Introduction}

Federated Learning (FL) introduces unique challenges that require adapted MLOps practices. Unlike centralized ML, FL involves:
\begin{itemize}
    \item Distributed training across multiple clients
    \item Privacy constraints (no access to client data)
    \item Round-based iterative training
    \item Model aggregation instead of direct training
    \item Communication overhead management
\end{itemize}

\section{Key Differences: Centralized ML vs Federated Learning}

\begin{table}[h]
\centering
\begin{tabular}{|l|l|l|}
\hline
\textbf{Aspect} & \textbf{Centralized ML} & \textbf{Federated Learning} \\
\hline
\textbf{Training} & Single location & Distributed across clients \\
\hline
\textbf{Data Access} & Full dataset available & Data stays on clients \\
\hline
\textbf{Model Updates} & Direct gradient updates & Weight aggregation (FedAvg) \\
\hline
\textbf{Monitoring} & Single training run & Multiple rounds + clients \\
\hline
\textbf{Versioning} & Model versions & Round versions + aggregated models \\
\hline
\textbf{Deployment} & Single model & Global model + client models \\
\hline
\end{tabular}
\caption{Centralized ML vs Federated Learning}
\end{table}

\section{MLOps Components Adapted for FL}

\subsection{1. Experiment Tracking (FL-Specific)}

\subsubsection{What to Track in FL:}

\textbf{Per-Round Metrics:}
\begin{itemize}
    \item Global model accuracy/loss (aggregated)
    \item Client participation rate
    \item Round completion time
    \item Communication overhead (bytes sent/received)
    \item Convergence rate across rounds
\end{itemize}

\textbf{Per-Client Metrics (Privacy-Preserving):}
\begin{itemize}
    \item Local training loss/accuracy (aggregated, not individual)
    \item Number of samples per client
    \item Training time per client
    \item Model update contribution (weighted by samples)
\end{itemize}

\textbf{In Your Code:}
\begin{lstlisting}[language=Python]
# flower_server_metrics.py tracks:
- Round number
- Aggregated accuracy/loss
- Bytes sent/received
- Parameters count
- Round duration
- Client participation
\end{lstlisting}

\subsubsection{MLOps Tools for FL Experiment Tracking:}

\textbf{MLflow Integration:}
\begin{lstlisting}[language=Python]
import mlflow

# Track FL experiment
with mlflow.start_run():
    mlflow.log_param("num_rounds", 5)
    mlflow.log_param("num_clients", 2)
    mlflow.log_param("fraction_fit", 1.0)
    
    # Track per-round metrics
    for round_num in range(1, 6):
        mlflow.log_metric("round_accuracy", accuracy, step=round_num)
        mlflow.log_metric("round_loss", loss, step=round_num)
        mlflow.log_metric("bytes_sent", bytes_sent, step=round_num)
        mlflow.log_metric("round_time", round_time, step=round_num)
    
    # Log final aggregated model
    mlflow.keras.log_model(aggregated_model, "model")
\end{lstlisting}

\subsection{2. Model Versioning (Round-Based)}

\subsubsection{FL Model Versioning Strategy:}

\textbf{Current Approach:}
\begin{itemize}
    \item Single model saved: `MLPv2_FL.h5`
    \item No versioning of rounds
    \item No tracking of intermediate rounds
\end{itemize}

\textbf{MLOps Approach:}
\[
\text{Model Version} = \text{Model\_Type}\_\text{Round}\_\text{Accuracy}\_\text{Timestamp}
\]

\textbf{Example:}
\begin{itemize}
    \item `MLPv2_FL_Round1_Acc0.85_20250115.h5`
    \item `MLPv2_FL_Round3_Acc0.91_20250115.h5`
    \item `MLPv2_FL_Round5_Acc0.94_20250115.h5` (final)
\end{itemize}

\textbf{Implementation:}
\begin{lstlisting}[language=Python]
# After each round aggregation
round_accuracy = aggregated_metrics.get('accuracy', 0)
model_version = f"{model_type}_FL_Round{round_num}_Acc{round_accuracy:.3f}_{timestamp}"
model.save(f"models/{model_version}.h5")

# Register in model registry
mlflow.register_model(
    f"runs:/{run_id}/model",
    f"{model_type}_FL_v{round_num}"
)
\end{lstlisting}

\subsubsection{Why Round-Based Versioning Matters:}

\begin{itemize}
    \item \textbf{Rollback capability:} Can revert to previous round if performance degrades
    \item \textbf{Convergence analysis:} Track how model improves across rounds
    \item \textbf{Debugging:} Identify which round introduced issues
    \item \textbf{Research:} Compare different FL strategies (FedAvg vs FedProx vs others)
\end{itemize}

\subsection{3. Model Registry (FL-Aware)}

\subsubsection{Registry Structure:}

\[
\text{Registry} = \begin{cases}
\text{Global Models} & \text{(Aggregated, server-side)} \\
\text{Client Models} & \text{(Local, client-side, optional)}
\end{cases}
\]

\textbf{Global Model Registry:}
\begin{itemize}
    \item \textbf{Entry:} Aggregated model after each round
    \item \textbf{Metadata:} Round number, accuracy, loss, client count, aggregation method
    \item \textbf{Staging:} Round-by-round progression
    \item \textbf{Production:} Final model after all rounds
\end{itemize}

\textbf{Example Registry Entry:}
\begin{lstlisting}[language=JSON]
{
  "model_id": "MLPv2_FL_Round5",
  "version": "5.0",
  "round": 5,
  "accuracy": 0.94,
  "loss": 0.12,
  "num_clients": 2,
  "total_samples": 72235,
  "aggregation_method": "FedAvg",
  "timestamp": "2025-01-15T10:30:00",
  "artifacts": {
    "model_path": "models/MLPv2_FL_Round5.h5",
    "metrics_path": "metrics/round5_metrics.json"
  }
}
\end{lstlisting}

\subsection{4. Monitoring (FL-Specific Metrics)}

\subsubsection{What to Monitor in FL:}

\textbf{Training Metrics:}
\begin{itemize}
    \item \textbf{Convergence:} How quickly global model converges
    \item \textbf{Stability:} Variance in client updates
    \item \textbf{Participation:} Client dropout rate, availability
    \item \textbf{Communication:} Network overhead, round time
\end{itemize}

\textbf{Model Quality Metrics:}
\begin{itemize}
    \item \textbf{Global accuracy:} Aggregated model performance
    \item \textbf{Client accuracy variance:} How consistent are client models
    \item \textbf{Data distribution:} IID vs non-IID detection
    \item \textbf{Model drift:} Performance degradation over rounds
\end{itemize}

\textbf{In Your Code:}
\begin{lstlisting}[language=Python]
# flower_server_metrics.py already tracks:
- Round-by-round accuracy/loss
- Bytes sent/received per round
- Round duration
- Client participation
- Parameter counts
\end{lstlisting}

\subsubsection{MLOps Monitoring Dashboard:}

\textbf{Key Metrics to Display:}
\[
\begin{aligned}
\text{Convergence Rate} &= \frac{\text{Accuracy}_{\text{round }n} - \text{Accuracy}_{\text{round }1}}{\text{round }n - 1} \\
\text{Communication Efficiency} &= \frac{\text{Model Improvement}}{\text{Bytes Transferred}} \\
\text{Client Contribution} &= \frac{n_k \cdot \text{Accuracy}_k}{\sum n_k}
\end{aligned}
\]

\subsection{5. CI/CD Pipeline (FL-Aware)}

\subsubsection{FL Pipeline Stages:}

\textbf{Stage 1: Data Validation (Client-Side)}
\begin{itemize}
    \item Validate data format on each client
    \item Check data quality (missing values, outliers)
    \item Verify data distribution (IID/non-IID)
    \item \textbf{Privacy:} Validation happens locally, only metadata sent
\end{itemize}

\textbf{Stage 2: Model Initialization (Server-Side)}
\begin{itemize}
    \item Initialize global model architecture
    \item Set hyperparameters (learning rate, epochs per round)
    \item Configure FL strategy (FedAvg parameters)
\end{itemize}

\textbf{Stage 3: FL Training Loop (Distributed)}
\begin{itemize}
    \item \textbf{For each round:}
    \begin{enumerate}
        \item Server sends global weights to clients
        \item Clients train locally (parallel)
        \item Clients send updated weights back
        \item Server aggregates weights (FedAvg)
        \item Server evaluates aggregated model
        \item Log metrics, version model
    \end{enumerate}
\end{itemize}

\textbf{Stage 4: Model Validation (Server-Side)}
\begin{itemize}
    \item Check convergence (accuracy improvement)
    \item Validate model performance (threshold check)
    \item Compare with previous rounds
    \item Decide: continue training or deploy
\end{itemize}

\textbf{Stage 5: Deployment (Server-Side)}
\begin{itemize}
    \item Deploy aggregated global model
    \item Update model registry
    \item Notify clients of new model version
    \item Monitor production performance
\end{itemize}

\subsubsection{GitHub Actions Example:}

\begin{lstlisting}[language=YAML]
name: FL Training Pipeline

on:
  schedule:
    - cron: '0 0 * * 0'  # Weekly retraining
  workflow_dispatch:

jobs:
  fl-training:
    runs-on: ubuntu-latest
    steps:
      - name: Validate Data
        run: python validate_data.py
        
      - name: Start FL Server
        run: |
          docker-compose up -d flower-server-mlpv2
          
      - name: Start FL Clients
        run: |
          docker-compose up -d flower-worker-1 flower-worker-2
          
      - name: Monitor FL Training
        run: |
          python monitor_fl_training.py --wait-for-completion
          
      - name: Evaluate Model
        run: |
          python evaluate_fl_model.py --model MLPv2_FL.h5
          
      - name: Register Model
        if: success()
        run: |
          mlflow models register --model-path models/MLPv2_FL.h5
          
      - name: Deploy if Improved
        if: success()
        run: |
          python deploy_if_improved.py
\end{lstlisting}

\subsection{6. Model Deployment (FL-Specific)}

\subsubsection{Deployment Architecture:}

\[
\text{FL Deployment} = \begin{cases}
\text{Global Model} & \text{(Server, aggregated)} \\
\text{Client Models} & \text{(Optional, local inference)}
\end{cases}
\]

\textbf{Option 1: Centralized Deployment (Current)}
\begin{itemize}
    \item Deploy aggregated global model on server
    \item Clients send inference requests to server
    \item \textbf{Pros:} Single model to maintain, consistent predictions
    \item \textbf{Cons:} Network latency, server load
\end{itemize}

\textbf{Option 2: Federated Deployment}
\begin{itemize}
    \item Deploy global model to each client
    \item Clients perform local inference
    \item \textbf{Pros:} Low latency, privacy-preserving
    \item \textbf{Cons:} Model synchronization needed
\end{itemize}

\textbf{Option 3: Hybrid}
\begin{itemize}
    \item Global model on server for centralized inference
    \item Client models for local inference (optional)
    \item \textbf{Pros:} Flexibility, fallback options
\end{itemize}

\subsubsection{Deployment Pipeline:}

\begin{lstlisting}[language=Python]
# After FL training completes
def deploy_fl_model(model_path, model_version):
    # 1. Load aggregated model
    global_model = load_model(model_path)
    
    # 2. Validate performance
    test_accuracy = evaluate_model(global_model)
    if test_accuracy < threshold:
        raise ValueError("Model performance below threshold")
    
    # 3. Register in model registry
    mlflow.register_model(model_path, f"MLPv2_FL_{model_version}")
    
    # 4. Deploy to production
    # Option A: Server deployment
    deploy_to_server(global_model)
    
    # Option B: Client deployment (federated)
    deploy_to_clients(global_model)
    
    # 5. Monitor deployment
    start_monitoring(global_model)
\end{lstlisting}

\subsection{7. Automated Retraining (FL-Specific)}

\subsubsection{Retraining Triggers:}

\textbf{1. Scheduled Retraining:}
\begin{itemize}
    \item Weekly/Monthly FL rounds
    \item New data accumulates on clients
    \item Global model needs updates
\end{itemize}

\textbf{2. Performance Degradation:}
\begin{itemize}
    \item Monitor production accuracy
    \item If accuracy drops below threshold → trigger FL retraining
    \item New attack patterns detected
\end{itemize}

\textbf{3. Data Drift Detection:}
\begin{itemize}
    \item Detect distribution shift in client data
    \item New DDoS attack types emerge
    \item Network behavior changes
\end{itemize}

\textbf{4. New Clients Join:}
\begin{itemize}
    \item New network nodes added
    \item Retrain to incorporate new client data
    \item Maintain model freshness
\end{itemize}

\subsubsection{FL Retraining Pipeline:}

\begin{lstlisting}[language=Python]
def fl_retraining_pipeline():
    # 1. Check if retraining needed
    if not should_retrain():
        return
    
    # 2. Load previous model as initialization
    previous_model = load_model("models/MLPv2_FL_latest.h5")
    
    # 3. Start FL training with warm start
    fl_server.start_training(
        initial_weights=previous_model.get_weights(),
        num_rounds=5
    )
    
    # 4. Monitor training
    monitor_fl_rounds()
    
    # 5. Compare with previous model
    new_accuracy = evaluate_new_model()
    old_accuracy = evaluate_old_model()
    
    # 6. Deploy if improved
    if new_accuracy > old_accuracy:
        deploy_new_model()
    else:
        keep_old_model()
\end{lstlisting}

\section{FL-Specific MLOps Challenges}

\subsection{1. Privacy Constraints}

\textbf{Challenge:} Cannot access client data for monitoring/validation

\textbf{MLOps Solution:}
\begin{itemize}
    \item \textbf{Federated metrics:} Aggregate metrics without exposing individual client data
    \item \textbf{Differential privacy:} Add noise to aggregated metrics
    \item \textbf{Secure aggregation:} Use cryptographic protocols
    \item \textbf{Metadata only:} Track only aggregated statistics
\end{itemize}

\textbf{In Your Code:}
\begin{lstlisting}[language=Python]
# Privacy-preserving metrics
aggregated_accuracy = weighted_average(client_accuracies, client_samples)
# No individual client data exposed
\end{lstlisting}

\subsection{2. Communication Overhead}

\textbf{Challenge:} Model weights are large, multiple rounds increase communication

\textbf{MLOps Solution:}
\begin{itemize}
    \item \textbf{Compression:} Quantization, pruning, compression
    \item \textbf{Selective updates:} Only send significant weight changes
    \item \textbf{Monitoring:} Track bytes sent/received (you already do this!)
    \item \textbf{Optimization:} Reduce model size, use efficient formats
\end{itemize}

\textbf{In Your Code:}
\begin{lstlisting}[language=Python]
# Already tracking communication
bytes_sent = self._estimate_bytes(weights_bytes)
params_sent = self._count_parameters(weights)
# Could optimize: compression, quantization
\end{lstlisting}

\subsection{3. Client Heterogeneity}

\textbf{Challenge:} Different clients have different:
\begin{itemize}
    \item Data distributions (IID vs non-IID)
    \item Data sizes
    \item Computational resources
    \item Network conditions
\end{itemize}

\textbf{MLOps Solution:}
\begin{itemize}
    \item \textbf{Adaptive strategies:} FedProx, FedNova for heterogeneous clients
    \item \textbf{Client selection:} Select clients based on resources/data
    \item \textbf{Monitoring:} Track client participation, performance variance
    \item \textbf{Fairness:} Ensure all clients benefit from FL
\end{itemize}

\subsection{4. Round-Based Training}

\textbf{Challenge:} Training happens in rounds, not continuous

\textbf{MLOps Solution:}
\begin{itemize}
    \item \textbf{Round tracking:} Version models per round
    \item \textbf{Convergence monitoring:} Track improvement across rounds
    \item \textbf{Early stopping:} Stop if no improvement
    \item \textbf{Checkpointing:} Save models after each round
\end{itemize}

\section{MLOps Implementation for Your FL System}

\subsection{Current State}

\textbf{What You Have:}
\begin{itemize}
    \item ✅ Basic monitoring (Flower dashboard)
    \item ✅ Model saving (H5 files)
    \item ✅ Metrics tracking (accuracy, loss, bytes)
    \item ✅ Docker containerization
    \item ✅ Round-based training
\end{itemize}

\textbf{What's Missing:}
\begin{itemize}
    \item ❌ Model versioning
    \item ❌ Experiment tracking (MLflow)
    \item ❌ Model registry
    \item ❌ Automated pipelines
    \item ❌ Production deployment
    \item ❌ Automated retraining
\end{itemize}

\subsection{Recommended MLOps Stack for FL}

\subsubsection{Tier 1: Essential (Start Here)}

\textbf{1. MLflow for FL Experiment Tracking}
\begin{lstlisting}[language=Python]
# Track FL experiments
import mlflow

mlflow.set_experiment("DDoS_Detection_FL")

with mlflow.start_run():
    # Log FL configuration
    mlflow.log_params({
        "num_rounds": 5,
        "num_clients": 2,
        "model_type": "MLPv2",
        "aggregation": "FedAvg"
    })
    
    # Track per-round metrics
    for round_num in range(1, 6):
        mlflow.log_metrics({
            "round_accuracy": accuracy[round_num],
            "round_loss": loss[round_num],
            "bytes_sent": bytes_sent[round_num],
            "round_time": round_time[round_num]
        }, step=round_num)
    
    # Log final model
    mlflow.keras.log_model(final_model, "model")
\end{lstlisting}

\textbf{2. Model Versioning}
\begin{lstlisting}[language=Python]
# Version models by round
def save_fl_model(model, round_num, accuracy, model_type):
    timestamp = datetime.now().strftime("%Y%m%d_%H%M%S")
    version = f"{model_type}_FL_Round{round_num}_Acc{accuracy:.3f}_{timestamp}"
    model_path = f"models/{version}.h5"
    model.save(model_path)
    
    # Also save as latest
    latest_path = f"models/{model_type}_FL_latest.h5"
    model.save(latest_path)
    
    return model_path
\end{lstlisting}

\subsubsection{Tier 2: Production (Next Steps)}

\textbf{3. FL-Aware CI/CD Pipeline}
\begin{itemize}
    \item Automated FL training on schedule
    \item Model validation after each round
    \item Automated deployment if improved
    \item Rollback capability
\end{itemize}

\textbf{4. Production Monitoring}
\begin{itemize}
    \item Monitor global model performance
    \item Track client participation
    \item Detect data drift
    \item Alert on performance degradation
\end{itemize}

\subsubsection{Tier 3: Advanced}

\textbf{5. Federated Model Serving}
\begin{itemize}
    \item Deploy models to clients
    \item Local inference capabilities
    \item Model synchronization
\end{itemize}

\textbf{6. Advanced FL Strategies}
\begin{itemize}
    \item FedProx (handles heterogeneity)
    \item Secure aggregation
    \item Differential privacy integration
\end{itemize}

\section{FL-Specific MLOps Metrics}

\subsection{Key Performance Indicators (KPIs)}

\textbf{Training Metrics:}
\[
\begin{aligned}
\text{Convergence Rate} &= \frac{\Delta \text{Accuracy}}{\text{Rounds}} \\
\text{Communication Efficiency} &= \frac{\text{Accuracy Gain}}{\text{Total Bytes}} \\
\text{Client Participation Rate} &= \frac{\text{Active Clients}}{\text{Total Clients}} \\
\text{Round Efficiency} &= \frac{\text{Accuracy Improvement}}{\text{Round Time}}
\end{aligned}
\]

\textbf{Model Quality Metrics:}
\[
\begin{aligned}
\text{Global Accuracy} &= \text{Aggregated test accuracy} \\
\text{Client Variance} &= \text{Var}(\text{Client Accuracies}) \\
\text{Model Stability} &= 1 - \frac{\text{Var}(\text{Round Accuracies})}{\text{Mean}(\text{Round Accuracies})}
\end{aligned}
\]

\textbf{System Metrics:}
\[
\begin{aligned}
\text{Communication Overhead} &= \frac{\text{Total Bytes}}{\text{Model Parameters}} \\
\text{Training Time} &= \sum \text{Round Times} \\
\text{Client Dropout Rate} &= \frac{\text{Clients Dropped}}{\text{Total Clients}}
\end{aligned}
\]

\section{Implementation Roadmap}

\subsection{Phase 1: Foundation (Weeks 1-2)}

\begin{enumerate}
    \item \textbf{Add MLflow:} Integrate experiment tracking
    \item \textbf{Model versioning:} Implement round-based versioning
    \item \textbf{Enhanced logging:} Structured logging for FL rounds
\end{enumerate}

\subsection{Phase 2: Automation (Weeks 3-4)}

\begin{enumerate}
    \item \textbf{CI/CD pipeline:} Automated FL training
    \item \textbf{Model registry:} Centralized model management
    \item \textbf{Automated evaluation:} Post-training validation
\end{enumerate}

\subsection{Phase 3: Production (Weeks 5-6)}

\begin{enumerate}
    \item \textbf{Model deployment:} API endpoints for inference
    \item \textbf{Production monitoring:} Real-time performance tracking
    \item \textbf{Automated retraining:} Trigger-based FL rounds
\end{enumerate}

\section{Conclusion}

MLOps is \textbf{highly applicable} to Federated Learning, but requires adaptations:

\begin{itemize}
    \item \textbf{Experiment tracking} → Track FL rounds, not just final model
    \item \textbf{Model versioning} → Version by round, track convergence
    \item \textbf{Monitoring} → FL-specific metrics (convergence, communication)
    \item \textbf{Deployment} → Consider federated vs centralized deployment
    \item \textbf{Privacy} → Privacy-preserving monitoring and metrics
\end{itemize}

\textbf{Key Insight:} MLOps for FL focuses on the \textbf{aggregation process} and \textbf{round-based training}, not just the final model. The entire FL lifecycle needs operationalization.

\end{document}

